% !TEX program = pdflatex
% Rebuilt main.tex on 2026-06-16 to force stale-copy replacement.
\documentclass[a4paper,twocolumn,10pt]{article}

\usepackage[top=30mm,bottom=27mm,left=18mm,right=18mm]{geometry}
\setlength{\columnsep}{7mm}
\usepackage[whole]{bxcjkjatype}
\usepackage{booktabs}
\usepackage{array}
\usepackage{graphicx}
\usepackage{setspace}
\usepackage{url}
\usepackage{cite}
\setstretch{1.05}
\pagestyle{empty}
\renewcommand{\tablename}{表}
\renewcommand{\figurename}{図}
\renewcommand{\refname}{参考文献}
\setlength{\textfloatsep}{7pt}
\setlength{\floatsep}{6pt}
\setlength{\intextsep}{6pt}
\setcounter{secnumdepth}{2}
\renewcommand{\thesection}{\arabic{section}}
\renewcommand{\thesubsection}{\thesection.\arabic{subsection}}

\newcommand{\code}[1]{{\small\ttfamily #1}}
\newcommand{\affmark}[1]{\footnotemark[#1]}

\begin{document}

\twocolumn[
\begin{center}
\vspace*{0mm}
{\large\bfseries
Windowsエンドポイントログを用いた正常行動の行動列復元による偽陽性判断支援の検討\par}
\vspace{1mm}
{\large
Reconstructing Benign Behavior Sequences from Windows Endpoint Logs for False Positive Triage\par}
\vspace{2mm}
{\small
\begin{tabular}{@{}cccc@{}}
小松崎　大世\affmark{1} &
横山　想一郎\affmark{2} &
山下　倫央\affmark{3} &
川村　秀憲\affmark{2} \\
Taisei Komatsuzaki &
Soichiro Yokoyama &
Tomohisa Yamashita &
Hidenori Kawamura
\end{tabular}
\par}
\end{center}
\vspace{3mm}
]
{\let\footnotesize\normalsize
\footnotetext[1]{北海道大学 大学院情報科学院, Graduate School of Information Science and Technology, Hokkaido University}
\footnotetext[2]{北海道大学 大学院情報科学研究院, Faculty of Information Science and Technology, Hokkaido University}
\footnotetext[3]{北海道大学 情報基盤センター, Information Initiative Center, Hokkaido University}}

\section{はじめに}

サイバー攻撃は，情報漏えい，業務停止，不正送金などを通じて組織活動に直接的な影響を与える．
そのため組織では，攻撃の疑いがある活動を早期に検知し，利用者端末やサーバ上で実際に何が起きたのかを確認する必要がある．
本稿では，このような調査対象となる端末やサーバをエンドポイントと呼ぶ．
エンドポイント検知・対応（Endpoint Detection and Response; EDR）は，エンドポイント上の活動を収集・分析し，攻撃の疑いを警告する仕組みである．
ただし，EDRの警告が出た時点で，その活動が攻撃なのか，管理作業やアプリケーション更新などの通常操作なのかは確定していない．
攻撃の見逃しは大きな被害につながるため，検知は疑わしい活動を広めに拾う設計になりやすく，個々の警告を確認するインシデント分析に負担が集まる．

インシデント分析では，調査者が警告文面だけでなくログに戻り，端末上の行動を確認する．
攻撃者は，後続の攻撃で再実行できる足掛かりを残すため，レジストリや自動起動設定などに値を登録する場合がある．
このような永続化に似た操作は，正規アプリケーションの更新や管理作業でも発生しうる．
そのため，「永続化が疑われる」という警告だけでは判断できず，どのプロセスが，どのレジストリパスに，どの値を設定したかを，対応するログとともに確認する必要がある．
この対応付けには，ログ種別，プロセス親子関係，組織内の通常操作に関する知識が必要であり，調査者の経験に依存しやすい．

インシデント分析の対象となるログは，本来，攻撃に由来する場合も通常操作に由来する場合もある．
EDRは攻撃の見逃しを避けるために疑わしい活動を広めに警告するため，正常行動を含むログも分析対象になりやすい．
本稿では，このような通常操作に由来する端末上の行動を正常行動と呼ぶ．
警告後の調査では，攻撃を見つけるだけでなく，正常行動をログ上の証跡と対応付けることも重要である．
一方，既存研究の多くは，異常検知や攻撃調査を主対象としており，警告後の正常行動を証跡付きで復元できるかは十分に検討されていない．

本稿では，警告が出たものの通常操作として説明できる場合を偽陽性と呼び，偽陽性判断支援のため，調査者に提示する証跡付き行動列をログから復元する方法を提案する．
偽陽性判断支援とは，警告を自動的に良性と判定することではなく，調査者が通常操作として説明できるかを確認するための行動列と証跡を提示することを指す．
偽陽性判断では，「攻撃らしい証拠が少ない」ことだけでなく，「通常操作としてどのように説明できるか」をログ上の証跡とともに示す必要がある．
このため，本研究では大規模言語モデル（Large Language Model; LLM）を用いたエージェントによる正常行動復元の可能性を検証する．

証跡付き行動列とは，どの主体が，どの対象に，どのような操作を，どの時刻または順序で行ったかを，ログ上の証跡とともに並べたものである．
主体はプロセスや利用者操作，対象はファイル，レジストリキー，通信先などに対応し，証跡とは各行動を確認するために参照できるログ行を指す．
本稿では，EDRの警告と，警告に含まれる関連プロセス名，ホスト，時刻などの情報を調査起点として用いる．
調査起点はログ集合と無関係な外部情報ではなく，ログ集合の中から調査すべき範囲を指定する手がかりである．
警告に含まれる要約や関連情報は探索範囲を定める手がかりとして用いるが，それ自体を行動の証跡とはみなさない．
したがって，本研究で復元すべきものは，警告文面の言い換えではなく，ログに基づいて調査者が確認できる行動列である．

対象とするログは，Windowsエンドポイントログである．
本稿がWindows端末に着目するのは，WindowsがデスクトップOSとして広く利用され，\(2026\)年\(5\)月時点で世界のデスクトップOS利用の\(62.16\%\)を占めるためである\cite{statcounter_desktop_os}．
本稿では，Windows環境で得られるプロセス親子関係，コマンドライン，レジストリ操作，ファイル操作，通信先などのログを用いて行動列復元を評価する．
なお，本稿の問題設定は，主体，操作，対象，時刻，証跡という抽象化に基づくため，同様の要素を確認できるログにも拡張可能な枠組みである．

本稿の貢献は次の三点である．
第一に，EDRの警告を受けた偽陽性判断支援を，証跡付き行動列復元の問題として定式化する．
第二に，調査起点，ログデータベース，ログ行，行動主体，対象，行動，証跡ログ，行動列の関係を記号により明確化する．
第三に，ATLASv2由来のWindowsエンドポイントログ\cite{atlasv2}を用い，警告要約への依存を切り分けながら，復元された行動列の品質を評価する．

\section{関連研究}

\subsection{インシデント分析の自動化}

ログ異常検知の研究では，ログ系列から通常と異なる振る舞いを検出する手法が提案されている．
DeepLog\cite{deeplog}は，システムログ系列の異常検知と診断を扱う代表的な研究である．
このような手法は，大量のログから注意すべき箇所を絞り込むうえで有用である．
一方，実際のインシデント分析では，検知された警告が真の攻撃かを確認し，関連する端末上の行動を追跡する必要がある．
NoDoze\cite{nodoze}は，警告に関係する因果依存グラフを作成し，過去の発生頻度に基づいて警告を優先付けすることで，警告疲れの軽減を目指す．
ATLAS\cite{atlas}は，監査ログから攻撃の手順を表す系列を構成し，攻撃ストーリーを復元する枠組みである．
AIRTAG\cite{airtag}は，因果グラフを明示的に作る代わりに，ログテキストを用いた教師なし学習により攻撃に関係する実体を推定する．
これらの研究は，検知後の調査を自動化する方向性を示しているが，主な対象は攻撃に関係する実体や攻撃ストーリーの復元である．

\subsection{攻撃知識ベースと評価データセット}

攻撃知識ベースも，エンドポイント調査の重要な支援情報である．
MITRE ATT\&CK\cite{attack}は，攻撃者の戦術や技術を整理した知識ベースである．
また，Living Off The Land Binaries, Scripts and Libraries（LOLBAS）\cite{lolbas}は，Windowsの正規バイナリやスクリプトが攻撃に悪用される事例を整理している．
これらは，ある操作が攻撃技術と似た外形を持つかを考える手がかりになる．
しかし，攻撃技術に似た操作であっても，実際には正規アプリケーションの更新や管理作業として発生する場合がある．
また，ATLASv2\cite{atlasv2}はATLASの攻撃シナリオを再現しつつ，SysmonログやCarbon Black Cloudを含む観測点と，人手による背景行動を追加したデータセットである．

\subsection{大規模言語モデルの登場とマルチエージェント調査支援}

大規模言語モデル（LLM）を用いた調査支援も研究されている．
LLMは，自然言語で与えられた警告やログ断片から調査仮説を作り，確認すべき項目を検索質問へ変換し，得られた観測結果を説明文として整理できる可能性がある．
ReAct\cite{react}は，推論と言語的な行動を組み合わせて外部情報を探索する枠組みである．
CLOUSEAU\cite{clouseau}は，攻撃調査を階層的な複数エージェントで進めるログ探索基盤である．
複数エージェント構成は，調査論点の生成，検索質問への分解，SQL検索，証跡整理といった役割を分けられる点で，長いログを一度に要約する方法とは異なる．
本研究では，CLOUSEAU型のログ探索を実験基盤として利用する．
ただし，本稿の新規性は，エージェント構造そのものではない．
本稿の差分は，攻撃調査を前提としたCLOUSEAU型ログ探索を，警告後の未確定なWindows端末行動について，証跡付き行動列を復元する問題へ再設定した点にある．
具体的には，警告名や検知理由を行動として扱わないプロンプト制約，主体・操作・対象・時刻・証跡ログを返す出力形式，警告要約への依存を切り分ける入力条件と評価設計を追加する．

既存研究は，ログ収集，異常検知，攻撃知識の整理，LLMによる調査支援をそれぞれ前進させてきた．
しかし，警告直後の攻撃か通常操作かが未確定な状態で，通常操作として説明できるかを検討するための証跡付き行動列を復元する問題は，十分に整理されていない．

\section{問題設定}
\label{sec:problem}

本章では，本研究で扱う入力，出力，および復元対象を定義する．
本稿の対象は，警告直後の未確定な状態で，通常操作として説明できる行動列とその証跡をそろえる問題である．
本研究の入力は，ログデータベース \(D\) と調査起点 \(p\) の組である．
ログデータベースは，調査対象となる端末活動の観測結果を検索可能にまとめたものであり，調査起点は，その中から確認すべき範囲や対象を指定する手がかりである．
本研究では警告文面だけを入力として扱うのではなく，EDRの警告に含まれる関連プロセス名，ホスト，時刻などの情報をログデータベース内の確認範囲を定める手がかりとして扱い，調査者が確認できる行動列を出力する．

\subsection{ログデータベースとログ行}

本稿では，端末上の活動を記録した複数種類のログを検索可能な形にまとめたものをログデータベースと呼ぶ．
ログデータベースを
\[
D=\{r_1,\ldots,r_n\}
\]
とする．
各ログ行 \(r_i\) は，
\[
r_i=(t_i,h_i,u_i,s_i,v_i,o_i,c_i,\ell_i)
\]
で表す．
ここで，\(t_i\) は時刻，\(h_i\) はホスト，\(u_i\) は利用者またはアカウント，\(s_i\) はプロセスや利用者操作などの行動主体，\(v_i\) はプロセス起動，ファイル作成，レジストリ設定，通信などの操作，\(o_i\) は実行ファイル，ファイルパス，レジストリキー，通信先などの対象，\(c_i\) はコマンドラインやレジストリ値などの補足情報，\(\ell_i\) はログ行識別子である．

\subsection{調査起点と調査対象ログ}

調査起点を
\[
p=(h_0,x_0,t_0,\Delta_-,\Delta_+,a_0)
\]
とする．
ここで，\(h_0\) は調査対象のホスト，\(x_0\) は調査の手がかりとなるプロセス，ファイル，レジストリ，通信先など，\(t_0\) は基準時刻である．
\(\Delta_-\) と \(\Delta_+\) は，基準時刻の前後で調査対象とする時間幅であり，EDRの警告が発生した時刻，調査者が確認したい操作の時刻，または実験条件に基づいて事前に与えられる．
\(a_0\) は，EDRの警告に含まれる任意の追加情報であり，存在しない場合もある．
実運用では，EDRが警告名，検知理由，関連プロセス，発生時刻を提示する場合と，調査者が画面表示，利用者申告，別システムの通知などからホスト，プロセス，時刻だけを手動で与える場合がある．
本問題設定では，前者を \(a_0\) を含む調査起点，後者を \(a_0\) を含まない調査起点として扱う．
時間窓を
\[
W(p)=[t_0-\Delta_-,t_0+\Delta_+]
\]
と定義する．
このとき，調査対象ログを
\[
D_p=\{r_i\in D \mid h_i=h_0,\ t_i\in W(p)\}
\]
とする．
ここで \(D_p\) は，ホストと時間窓で切り出した候補ログ集合である．
\(x_0\) や \(a_0\) は，\(D_p\) 内で関連するログ行や接続関係を示す手がかりである．
すなわち，調査範囲を一点に限定せず，ホストと時間窓で候補範囲 \(D_p\) を定める．

\subsection{出力と復元目標}

本研究の出力は，自然文の結論だけではなく，主体，行動，対象，時刻または順序，証跡ログを持つ行動列である．
行動ステップを
\[
y_j=(\tau_j,s_j,v_j,o_j,E_j)
\]
と定義する．
ここで，\(\tau_j\) は時刻または行動列内の順序，\(s_j\) は主体，\(v_j\) は行動，\(o_j\) は対象，\(E_j\subseteq D_p\) は当該行動を支える証跡ログ集合である．
本稿では，\(\tau_j\)，\(s_j\)，\(v_j\)，\(o_j\) を行動部品と呼ぶ．
行動部品は，行動列として復元すべき時刻または順序，主体，操作，対象である．
また，\(E_j\) のうち，当該行動主張を確認するために最低限必要なログ根拠を重要証跡と呼ぶ．
重要証跡は，単に時刻が近い周辺ログではなく，主体，操作，対象の対応を確認するために必要なログ行またはログ項目である．
偽陽性判断支援では，行動部品が復元されるだけでは不十分であり，調査者がその行動を通常操作として説明できるかを確認するための重要証跡と，行動の前後関係を示す順序が必要である．
また，正解行動列に属さない主張が多い場合，調査者が追加で確認すべき候補が増えるため，過剰主張を少なくすることも支援の条件となる．
\(E_j\) は空であってはならず，\(s_j\)，\(v_j\)，\(o_j\) は \(E_j\) 内のログ行またはログ行間の接続関係で確認できる必要がある．
主体，行動，対象は一つのログ行に揃う必要はなく，プロセス親子関係，コマンドライン，レジストリパス，通信先などを同一ホスト，近接時刻，同一プロセス系列で接続したまとまりも一つの行動ステップの証跡とする．
一方，時刻が近いだけで当該行動に寄与しないログは \(E_j\) に含めず，別の行動ステップまたは主系列外の事象として扱う．
復元すべき行動列を
\[
Y=(y_1,\ldots,y_m)
\]
とする．
\(Y\) は，時刻または因果的な接続に基づいて並べられる．
本問題の目的は，入力 \(D\) と \(p\) から，警告名や検知理由の言い換えではなく，ログ行に基づく推定行動列 \(\hat{Y}\) を出力することである．
例えば，EDRの警告に含まれる情報 \(a_0\) に「不審なレジストリ操作」と記載されていても，出力では主体，レジストリキー，値，時刻を証跡ログ \(E_j\) とともに示す必要がある．
評価においては，対象シナリオでログ上確認可能な行動列を正解行動列 \(Y^\ast\) と呼ぶ．
\(Y^\ast\) はモデル出力から作るものではなく，評価対象シナリオのログに基づいて事前に定める参照行動列である．
望ましい \(\hat{Y}\) は，\(Y^\ast\) に含まれる主体，行動，対象，順序，証跡を多く含み，かつ \(Y^\ast\) に属さない主張を少なくするものである．

以上の問題設定に基づき，次章では，調査起点からログを探索し，証跡付き行動列を復元する手法を述べる．

\section{提案手法}

\subsection{システム構成}

本研究で用いるシステムは，CLOUSEAU型のエージェント型ログ探索\cite{clouseau}を基盤とし，問題設定で定義した調査対象ログ \(D_p\) から証跡付き行動列 \(\hat{Y}\) を構成する．
ここでエージェントとは，調査論点の生成や調査実行を担うLLMモジュール，または検索質問をSQLとして実行する検索モジュールを指す．
本システムの目的は，調査者が通常操作として説明できるかを確認するための証跡付き行動列 \(\hat{Y}\) を提示することである．
警告を攻撃または通常操作に自動判定するものではない．

本システムは，調査対象ログ \(D_p\) を一度にすべて要約しない．
まず，調査起点から確認すべき論点を作り，その論点を具体的な検索質問へ分解する．
次に，SQL検索で該当ログ行を取得し，得られたログをもとに次の検索質問を作る．
この反復で行動の根拠になりうるログを集め，最後にそれらを主体，操作，対象，時刻，証跡に整理して行動列 \(\hat{Y}\) を構成する．

処理は，論点生成・出力構成，調査実行，SQL検索の三つに分かれる．
論点生成・出力構成はChief Agentが担い，調査起点 \(p\) と問題設定で定義した調査対象ログ \(D_p\) から，確認すべき調査論点を作る．
調査実行はInvestigator Agentが担い，調査論点と既存の検索結果から確認項目を具体化し，検索質問を作る．
SQL検索はSQL Agentが担い，検索質問に基づいてSQLite化されたログデータベースを検索し，ログ行集合を返す．
調査実行とSQL検索を必要に応じて反復し，得られた検索結果をInvestigator Agentが整理してChief Agentに戻す．
Chief Agentは，集められたログと調査論点を照合し，行動列 \(\hat{Y}\) を構成する．

以下では，この処理を記号で表す．
第 \(k\) 回目の検索質問を \(q_k\) とし，SQL検索で得られるログ行集合を
\[
R_k=Q(D_p,q_k)\subseteq D_p
\]
と表す．
ここで，\(Q\) は時刻範囲，ホスト，プロセス識別子，親子関係，ファイルパス，レジストリキー，通信先などに基づいてログを問い合わせる検索操作であり，本実装ではSQL Agentが担う．
検索質問 \(q_k\) は，調査起点 \(p\) と以前の検索結果 \(R_1,\ldots,R_{k-1}\) に基づいて決められる．
この反復により得られた観測ログ集合を
\[
O_K=\bigcup_{k=1}^{K}R_k
\]
とする．
本稿では，ログ探索を \(O_K=S(D,p)\)，行動列構成を \(\hat{Y}=G(O_K,p)\) と表す．
この流れにより，警告要約の言い換えではなく，SQL検索で得た証跡に基づく正常行動列を提示する．
図\ref{fig:pipeline}に，入力，内部処理，出力の関係と，SQL検索が調査対象ログデータベースへ問い合わせる関係を示す．
CLOUSEAU型ログ探索から本稿で変更した点を表\ref{tab:clouseau_diff}に示す．

\begin{table*}[t]
\centering
\scriptsize
\caption{CLOUSEAU型ログ探索から本稿で変更した点}
\label{tab:clouseau_diff}
\setlength{\tabcolsep}{3pt}
\begin{tabular}{p{0.16\textwidth}p{0.35\textwidth}p{0.40\textwidth}}
\toprule
観点 & CLOUSEAU型ログ探索 & 本稿での設定 \\
\midrule
目的 & 攻撃調査を進め，攻撃に関係する証跡を整理する & 未確定な警告について，正常行動として説明可能かを検討する証跡付き行動列を復元する \\
入力 & 警告や攻撃調査の手がかりを起点にログを探索する & ホスト，プロセス，時刻を基本起点とし，警告要約を含む場合，入力に含めない場合，検索対象から除く場合を分ける \\
プロンプト & 攻撃調査の仮説生成，調査実行，SQL検索，結果整理を役割分担する & 警告名や検知理由を行動として扱わず，主体，操作，対象，時刻，証跡ログに基づく出力を要求する \\
評価 & 調査結果として攻撃関連情報を整理できるかを見る & 行動部品再現率，重要証跡再現率，順序，過剰主張，費用，実行時間を測定する \\
\bottomrule
\end{tabular}
\end{table*}

\begin{figure*}[t]
\centering
\scriptsize
\setlength{\tabcolsep}{3pt}
\renewcommand{\arraystretch}{1.05}
\begin{tabular}{@{}c@{\;}c@{\;}c@{\;}c@{\;}c@{}}
\fbox{\begin{minipage}{0.16\textwidth}\centering
\textbf{入力}\\
\vspace{0.5mm}
ログデータベース \(D\)\\
調査起点 \(p\)
\end{minipage}}
&
\(\stackrel{p,D_p}{\longrightarrow}\)
&
\begin{minipage}{0.54\textwidth}\centering
\fbox{\begin{minipage}{0.95\linewidth}\centering
\textbf{内部処理}\\[-0.2mm]
\begin{tabular}{@{}c@{}}
\fbox{\begin{minipage}{0.68\linewidth}\centering 論点生成・出力構成 \end{minipage}}\\[-0.2mm]
\(\downarrow\) 調査論点 \quad \(\uparrow\) 観測ログ集合 \(O_K\)\\[-0.2mm]
\fbox{\begin{minipage}{0.68\linewidth}\centering 調査実行 \end{minipage}}\\[-0.2mm]
\(\downarrow\) 検索質問 \(q_k\) \quad \(\uparrow\) SQL検索結果 \(R_k\)\\[-0.2mm]
\begin{tabular}{@{}c@{\;}c@{\;}c@{}}
\fbox{\begin{minipage}{0.28\linewidth}\centering SQL検索 \end{minipage}}
&
\begin{tabular}{@{}c@{}}
\(\longrightarrow\)\\[-0.8mm]
{\scriptsize SQL問い合わせ}\\[-0.2mm]
\(\longleftarrow\)\\[-0.8mm]
{\scriptsize 該当ログ行}
\end{tabular}
&
\fbox{\begin{minipage}{0.34\linewidth}\centering 調査対象ログ\\データベース \(D_p\) \end{minipage}}
\end{tabular}\\[-0.2mm]
\multicolumn{1}{c}{\scriptsize 反復により \(O_K=\bigcup_{k=1}^{K}R_k\) を作り，論点生成・出力構成へ戻す}
\end{tabular}
\end{minipage}}
\end{minipage}
&
\(\stackrel{\hat{Y}}{\longrightarrow}\)
&
\fbox{\begin{minipage}{0.16\textwidth}\centering
\textbf{出力}\\
\vspace{0.5mm}
\(\hat{Y}\)\\
証跡付き行動列
\end{minipage}}
\end{tabular}
\caption{入力，内部処理，出力の関係とSQL検索におけるデータベース入出力}
\label{fig:pipeline}
\end{figure*}

\subsection{プロンプトと出力形式}

本研究では，CLOUSEAUの階層構造とSQL検索を含む役割分担を残しつつ，LLMを用いる各処理のプロンプトを本問題向けに差し替えた．
プロンプトは，役割，入力情報，制約条件，参照可能な観測結果，出力形式を明示する形に統一した．
論点生成では，警告名，検知理由，攻撃分類を行動そのものとして扱わず，ホスト，プロセス，時刻を探索範囲を定める手がかりとして扱い，確認すべき調査論点を作るよう指示した．
調査実行では，未観測の親プロセス，コマンドライン，対象を仮定せず，調査論点と直前の検索結果に含まれる観測値から次の確認項目と検索質問を作るよう指示した．
SQL検索では，SQL AgentがSQLiteを問い合わせ，元ログ行識別子，時刻，\code{PID/PPID}，プロセス名，コマンドライン，対象，ログ種別を保持して返すようにした．
出力構成では，SQL検索で得た観測ログ集合に基づき，根拠のない推定，警告文面の言い換え，主系列外の混入を避け，行動ステップごとに主体，操作，対象，時刻，根拠ログを含む形式で出力させる．
また，主系列に属さない近傍行動は，主行動ではなく除外または補助的な証跡として分けるよう指示した．
完全な実行プロンプトは実験コードとして保存し，補足資料として公開する．本文では追試に必要な主要な制約と出力形式を記述する．

\section{データセット}
\label{sec:dataset}

\subsection{ATLASv2}

本研究では，正常行動の復元を評価するため，攻撃シナリオだけでなく人手による背景行動を含むATLASv2\cite{atlasv2}を採用した．
ATLASv2は，ATLASの攻撃シナリオを再現した環境で収集された公開データセットであり，攻撃に関係するログに加えて，利用者操作やアプリケーション動作に由来する背景行動のログを含む．
対象ログには，Carbon Black Cloud（CBC）由来の端末活動イベント，Windows Security，Sysmonログ\cite{sysmon}，DNS，ブラウザ履歴などが含まれ，本研究ではこれらを検索可能なデータベースに変換した．
CBC由来ログは，検知器側の解釈を含む警告要約行と，プロセス実行，ファイル，レジストリ，通信などの端末活動イベントに分けて扱う．
警告要約行は調査起点または探索補助であり，行動そのものの正解証跡とはしない．

\subsection{評価対象行動列の選定}

本稿では，ATLASv2の背景行動のうち，正常行動として説明可能なログ範囲から\(23\)行動列を構成した．
評価対象の選定では，調査起点としてホスト，プロセス，時刻を与えられること，警告要約を除いてもログ上の証跡から正常行動として説明できること，復元難度の異なる行動を含むことを基準とした．
評価対象の選定とログ項目の解釈では，EDR由来ログの読み方，実務上の調査起点，ユースケースの妥当性について共同研究先の実務知見を参考にした．
正解行動列は，警告要約の記述をそのまま正解にせず，プロセス，コマンドライン，レジストリパス，通信証跡など，ログ上で確認可能な証跡に基づいて作成した．
作成基準はモデル出力を参照する前に固定し，警告名や検知理由だけに依存する行動は含めない．
また，警告要約行を用いなくても，残存する観測ログから回答可能な範囲に統一した．
本稿では，一つの行動列を一つの評価単位として扱い，単一または近接ログで確認しやすい明示的実行チェーン（\code{explicit}，\(15\)件），複数ログ系列を接続する複数段階チェーン（\code{multi-step}，\(7\)件），攻撃技術に似る外形を正常アプリケーション挙動として説明する意味判定チェーン（\code{semantic}，\(1\)件）に分類する．
EDRの警告に現れる正常行動には，単一の実行系列として確認しやすいものだけでなく，複数ログを接続しなければ説明できないものや，攻撃技術に似た外形を通常操作として意味付ける必要があるものが混在する．
この分類は，こうした行動列の性質の違いによって復元難度がどう変わるかを調べるために用いる．

\section{実験設計}
\label{sec:experiment}

\subsection{実験目的}

本実験の目的は，LLMエージェントが，偽陽性判断支援に必要な正常行動の証跡付き行動列をどの条件で復元できるかを明らかにすることである．
具体的には，次の四点を問う．
\begin{enumerate}
\setlength{\topsep}{2pt}
\setlength{\partopsep}{0pt}
\setlength{\itemsep}{0pt}
\setlength{\parsep}{0pt}
\item 調査起点から，行動部品，重要証跡，順序をどの程度復元できるか．これは，通常操作として説明できるかを調査者が判断するための最小単位である．
\item 警告要約の有無で復元品質がどう変化するか．実運用では，EDR画面の警告要約を使える場合と，ホスト，プロセス，時刻だけで調査を始める場合がある．
\item 内部モデルの違いが，復元品質，過剰主張，費用，実行時間にどう影響するか．SOCでは，多数の警告を扱うため，精度だけでなく確認負荷と運用負荷も制約になる．
\item 行動列の種類によって復元の難しさがどう異なるか．第\ref{sec:dataset}章で分類した行動列の種類ごとに，提案方法がどこまで復元条件を満たせるかを確認する．
\end{enumerate}

\subsection{実験設定}

本実験では，第\ref{sec:dataset}章で示した\(23\)行動列を用い，表\ref{tab:stage_setting}の三つの入力条件で評価する．
これらの入力条件は，第\ref{sec:problem}章で述べた調査起点の与えられ方の違いを切り分けるためのものである．
Stage \(1\)は，EDR画面や通知に警告要約，検知理由，関連プロセスが表示され，調査者がそれを初期手がかりとして使える状況を想定する．
Stage \(2\)は，調査者がホスト，プロセス，時刻だけを手動で与え，初期入力には警告要約を含めない状況を想定する．
ただし，検索対象のデータベース内には警告要約行が残っているため，ホストと時刻範囲で検索する過程では警告要約行にも到達しうる．
Stage \(3\)は，警告要約行を検索対象から除外し，EDR由来の端末活動イベントやSysmon等の観測ログから行動列を説明できるかを調べる条件である．
本実験では，各行動列について調査起点を中心とする\(5\)分幅の時間窓 \(W(p)\)，すなわち \(\Delta_-=\Delta_+=2.5\)分を事前に固定し，全Stage，全モデルで同一の時間窓を用いた．
これは，入力条件やモデル差の比較において，検索対象となる時間範囲の違いが結果に混入しないようにするためである．
警告要約には検知器側の解釈が含まれるため，Stage \(3\)ではCBCの警告要約行を除外し，CBC由来の端末活動イベント，Windows Security，Sysmonログ，DNS，ブラウザ履歴は残す．

\begin{table}[t]
\centering
\footnotesize
\caption{入力条件（Stage）の設計}
\label{tab:stage_setting}
\setlength{\tabcolsep}{3pt}
\begin{tabular}{p{0.16\linewidth}p{0.31\linewidth}p{0.41\linewidth}}
\toprule
Stage & 入力 & 意味 \\
\midrule
Stage 1 & CBC警告要約 + プロセス/時刻 & EDR画面の警告要約を初期手がかりとして使う状況 \\
Stage 2 & プロセス/時刻のみ & 手動起点だが検索によりDB内の警告要約へ到達可能 \\
Stage 3 & プロセス/時刻のみ & 警告要約に頼らず観測ログから復元 \\
\bottomrule
\end{tabular}
\end{table}

比較対象は\code{gpt-4.1-mini}，\code{gpt-5.4-mini}，および\code{gpt-5.5 low}である．
\code{gpt-4.1-mini}は軽量・低コスト側の基準，\code{gpt-5.4-mini}は本稿で主に提示候補として評価するモデル，\code{gpt-5.5 low}は復元能力の上限をみるためのプローブとして用いる．
\code{gpt-5.5 low}のlowは推論強度を低く設定した条件を表し，実運用で常に用いる提示用モデルではない．
表では紙面の都合上，それぞれ\code{4.1-mini}，\code{5.4-mini}，\code{5.5 low}と略記する．
各モデルについて，同一の\(23\)行動列，三つの入力条件，時間窓，評価指標で三回の評価を行い，各モデル\(207\)回の実行結果を得た．
Stage別の集計では，1モデル・1Stageあたり\(23\)行動列\(\times 3\)反復の\(69\)実行を用いる．
モデル別の全体集計では，各モデル\(207\)実行を用いる．
各実行について，入出力トークン数，利用費用，実行時間も記録した．
費用は記録済みの入出力トークン数と実行時点の価格設定に基づいて算出し，実行時間は調査起点を受け取って最終出力を返すまでの時間とする．

\subsection{評価方法}

評価では，第\ref{sec:problem}章で定義した行動部品と重要証跡に基づき，正解行動列 \(Y^\ast\) の各ステップとモデル出力の主張を照合する．
指標は，行動部品再現率，重要証跡再現率，候補精度，順序指標，過剰主張数である．
順序指標は，隣接する正解行動ステップのうち候補出力で順序関係が保持された数/\(\max(n_{\mathrm{gold}}-1,0)\)として算出する．
過剰主張は，証跡ログで支えられない主張，正解行動列 \(Y^\ast\) の主系列に属さない近傍事象，または警告名や検知理由の言い換えにとどまる主張として数える．
採点では，第一段階で候補出力の主張を行動部品と重要証跡の単位に分け，第二段階で正解行動列との対応と過剰主張を確認した．
表現ゆれだけでは不正解にしない一方，ログ上の証跡に対応しない主張は正解に含めない．
表では，行動部品再現率をAct.，重要証跡再現率をEv.，順序指標をOrd.，候補精度をPrec.，過剰主張総数をOver，過剰主張の一実行平均をO/r，入出力合計トークン数をTok.，実行時間中央値をp50と略記する．
ただし，本稿の採点は独立した複数評価者による一致度検証を含むものではなく，限られた評価者による二段階レビューである．

\section{実験結果}

\subsection{警告要約を除いた条件での復元}

表\ref{tab:stage_result}に入力条件別の結果を示す．
\code{gpt-5.4-mini}はStage \(1\)からStage \(3\)まで行動部品再現率が\(0.791\)，\(0.800\)，\(0.802\)とほぼ一定であり，重要証跡再現率もStage \(3\)で\(0.779\)を示した．
\code{gpt-5.5 low}でもStage \(3\)の行動部品再現率と重要証跡再現率はいずれも\(0.903\)であり，警告要約を除いても残存する観測ログから高い再現率で復元できる条件があることを示す．
ここで復元できる場合とは，警告要約を除いても，CBC由来の端末活動イベントやSysmon等に主体，操作，対象を結ぶ証跡が残る行動列を指す．
\code{gpt-4.1-mini}ではStage \(3\)の行動部品再現率が\(0.345\)，重要証跡再現率が\(0.164\)にとどまり，警告要約なしでの復元は内部モデルの能力に依存した．

\begin{table}[t]
\centering
\scriptsize
\caption{Stage別の復元性能}
\label{tab:stage_result}
\setlength{\tabcolsep}{0pt}
\begin{tabular*}{\linewidth}{@{\extracolsep{\fill}}llrrrrr@{}}
\toprule
Model & St. & Act. & Ev. & Ord. & Over & O/r \\
\midrule
\code{4.1-mini} & S1 & 0.639 & 0.292 & 0.341 & 422 & 6.12 \\
\code{4.1-mini} & S2 & 0.414 & 0.128 & 0.111 & 506 & 7.33 \\
\code{4.1-mini} & S3 & 0.345 & 0.164 & 0.151 & 262 & 3.80 \\
\code{5.4-mini} & S1 & 0.791 & 0.574 & 0.579 & 215 & 3.12 \\
\code{5.4-mini} & S2 & 0.800 & 0.754 & 0.540 & 213 & 3.09 \\
\code{5.4-mini} & S3 & 0.802 & 0.779 & 0.540 & 223 & 3.23 \\
\midrule
\code{5.5 low} & S1 & 0.937 & 0.908 & 0.897 & 374 & 5.42 \\
\code{5.5 low} & S2 & 0.915 & 0.908 & 0.857 & 415 & 6.01 \\
\code{5.5 low} & S3 & 0.903 & 0.903 & 0.849 & 446 & 6.46 \\
\bottomrule
\end{tabular*}
\par\raggedright\scriptsize 注: 各行は\(69\)実行の集計である．Overは過剰主張総数，O/rは一実行平均である．
\end{table}

\subsection{モデル別の復元品質}

表\ref{tab:overall}にモデル別の復元品質，過剰主張総数，費用，実行時間を示す．
\code{gpt-5.5 low}は行動部品再現率，重要証跡再現率，順序指標のいずれも最も高い一方，過剰主張総数，費用，実行時間も最大であった．
\code{gpt-5.4-mini}は\code{gpt-4.1-mini}に対して，行動部品再現率を\(0.466\)から\(0.798\)，重要証跡再現率を\(0.195\)から\(0.703\)へ改善し，過剰主張総数を\(1190\)から\(651\)へ減らした．
したがって，\code{gpt-4.1-mini}は軽量側の基準，\code{gpt-5.4-mini}は復元品質と確認負荷のバランスが良い提示候補，\code{gpt-5.5 low}は能力上限をみるプローブと位置付けられる．

\begin{table}[t]
\centering
\scriptsize
\caption{モデル別の復元品質，過剰主張，費用，実行時間}
\label{tab:overall}
\setlength{\tabcolsep}{0pt}
\begin{tabular*}{\linewidth}{@{\extracolsep{\fill}}lrrrrrrrrr@{}}
\toprule
Model & Act. & Ev. & Ord. & Prec. & Over & O/r & Tok. & USD & p50 \\
\midrule
\code{4.1-mini} & 0.466 & 0.195 & 0.201 & 0.366 & 1190 & 5.75 & 35.7 & 0.0176 & 4.00 \\
\code{5.4-mini} & 0.798 & 0.703 & 0.553 & 0.584 & 651 & 3.14 & 12.9 & 0.0210 & 1.37 \\
\midrule
\code{5.5 low} & 0.918 & 0.906 & 0.868 & 0.640 & 1235 & 5.97 & 37.1 & 0.4492 & 8.29 \\
\bottomrule
\end{tabular*}
\par\raggedright\scriptsize 注: 各値は\(207\)実行の集計である．Tok.は入出力合計トークン数，USDは費用，p50は実行時間中央値を表す．
\end{table}

過剰主張率は，\code{gpt-4.1-mini}が\(0.634\)，\code{gpt-5.4-mini}が\(0.416\)，\code{gpt-5.5 low}が\(0.273\)であり，\code{gpt-5.5 low}が最悪ではない．
本稿の結果は，強いモデルほど一主張あたり不正確という意味ではなく，多くの候補を出すため確認すべき過剰主張の絶対数が増えうることを示す．

\subsection{行動列種類別の復元品質}

復元能力と過剰主張抑制のバランスが良かった\code{gpt-5.4-mini}で見ると，明示的実行チェーンは行動部品再現率\(0.869\)，重要証跡再現率\(0.826\)，順序指標\(0.778\)，過剰主張\(2.56\)/runであった．
一方，複数段階チェーン（\code{multi-step}）では，それぞれ\(0.783\)，\(0.600\)，\(0.543\)，\(4.57\)/runとなり，複数ログを接続する行動列で復元が難しくなった．
意味判定チェーン（\code{semantic}）はDiscord Run Keyの一行動列のみであるため，考察で事例として扱う．

\section{考察}

本章では，第\ref{sec:dataset}章と第\ref{sec:experiment}章で設定した入力条件，モデル差，行動列分類に沿って，実験目的で示した四つの問いへの含意を整理する．

\subsection{警告要約に依存しない復元の意味}

第一の問いは，SOCで想定される調査起点から正常行動の証跡付き行動列をどの程度復元できるかであり，第二の問いは，その復元が警告要約にどの程度依存するかである．
Stage \(3\)で警告要約を除いても復元品質が大きく崩れなかったことは，警告文面の言い換えではなく，残存する観測ログから観測事実へ戻る経路を作れる可能性を示す．
これは，実務の初期調査で重要となる「警告文面を読むだけでなく，ログ上で何が起きたかを確認する」作業を支援するうえで意味がある．
ただし，Stage \(3\)でもEDR由来のイベント情報は残るため，本結果は任意の低レベルログのみで十分という主張ではない．

\subsection{モデル依存性と確認負荷}

第三の問いは，内部モデルの違いが復元品質と確認負荷にどう影響するかである．
本稿のモデル比較で重要なのは，復元能力と確認負荷が分離しうる点である．
\code{gpt-5.5 low}は能力上限を示す一方で，過剰主張総数，費用，実行時間も大きく，すべての警告に対する初期提示用モデルとしては重い．
\code{gpt-5.4-mini}は復元品質を大きく改善しつつ過剰主張総数を抑えたため，本実験では提示候補として最もバランスが良かった．
実運用では，軽量モデルで候補を提示し，証跡不足や複数段階チェーンで接続が難しい場合に強いモデルを用いる段階的運用が考えられる．

\subsection{費用と実行時間}

表\ref{tab:overall}では，\code{gpt-5.4-mini}が一実行平均\(0.0210\)米ドル，実行時間中央値\(1.37\)分で，\code{gpt-4.1-mini}より高い復元品質を短時間で示した．
一方，\code{gpt-5.5 low}は一実行平均\(0.4492\)米ドル，実行時間中央値\(8.29\)分であり，復元品質は高いが，多数の警告を順次処理する初期トリアージにはそのまま適用しにくい．
SOCでは同時に多数の警告を扱うため，費用と待ち時間は精度と同じく評価すべき制約である．
今後は，調査範囲の絞り込み，検索結果の再利用，段階的なモデル切替により，一件あたりの費用と待ち時間を抑えられるかを検証する．

\subsection{行動列の種類による難しさ}

第四の問いは，行動列の種類によって復元の難しさがどう異なるかである．
明示的実行チェーンでは主体，操作，対象，証跡の対応をログ上で取りやすい一方，複数段階チェーンでは時刻差，\code{PID/PPID}，同一利用者，同一ホスト，ファイルパス，通信先で複数ログを接続する必要がある．
難所は，「何が起きたか」を拾うことだけでなく，「どの事象を同じ主系列として接続してよいか」を判断する点にある．

\subsection{限界と今後の課題}

本稿の評価対象は公開評価データセット由来の\(23\)行動列であり，すべての企業環境，端末検知・対応製品，業務アプリケーション，運用ルールを代表するものではない．
また，本稿の実験はWindowsエンドポイントログを対象としており，他のOSやログ収集基盤での評価は行っていない．
実環境では，正解行動列が事前に存在しない場合が多く，調査者は不完全な警告，欠損したログ，組織固有の通常操作を前提に判断する必要がある．
正解行動列の作成，行動部品・重要証跡への分割，採点は証跡ベースで行い，作成基準をモデル出力参照前に固定したが，独立した複数評価者による一致度は報告していない．
ログ項目の読み方やユースケース選定では実務知見を参考にしたが，採点の独立性とは別である．
費用と実行時間は，本実験環境と実行時点の価格設定に依存するため，異なる価格設定，並列実行条件，検索対象ログ量では変化しうる．
今後は，複数評価者で正解行動列の作成と採点を分離し，\(\kappa\)やAC1などの一致度を報告する．
さらに，警告要約を含む起点と手動起点を実運用に近い形で分け，復元結果が判断時間，追加確認ログ件数，判断一貫性に与える影響を測定する．
また，各行動ステップに元ログ行識別子や具体的なログ項目を要求し，プロセス識別子，親プロセス識別子，時刻差，ファイルパスなどによる再ランキングで主系列を絞り込めるかを検証する．
これらは，復元精度だけでなく，実世界の初期調査で使える提示量，待ち時間，費用に収まるかを確認する次の実験課題である．

\section{おわりに}

本稿では，偽陽性判断支援を正常行動の証跡付き行動列復元として評価する枠組みを定め，ATLASv2由来の\(23\)行動列\cite{atlasv2}で検証した．
\code{gpt-5.4-mini}では，警告要約を除いたStage \(3\)でも行動部品再現率が\(0.802\)，重要証跡再現率が\(0.779\)となり，残存する観測ログから復元できる場合があることを示した．
一方，\code{gpt-5.5 low}は最も高い再現率を示すが過剰主張総数も最大であり，モデルの復元能力と確認負荷の小ささは一致しなかった．
また，明示的実行チェーンと複数段階チェーン（\code{multi-step}）では復元品質に差があり，行動列の種類によって難しさが異なることを示した．
今後は，復元結果を用いた調査者の判断時間や判断一貫性への効果を評価する．

\section*{謝辞}

本研究は，内閣府「地方大学・地域産業創生交付金」に採択された「次世代半導体をトリガーとした半導体の複合拠点の実現と地域経済の活性化」事業の支援を受けて実施しました．
本研究の実験にあたり，資料のご提供，ログ項目の解釈，実務上の調査状況，ユースケース選定に関するご助言をいただきました株式会社網屋（網屋さっぽろ研究所）の皆様に，心より感謝申し上げます．

\begin{thebibliography}{99}
\bibitem{sysmon}
Microsoft, ``Sysmon - Sysinternals,'' \url{https://learn.microsoft.com/en-us/sysinternals/downloads/sysmon} (accessed June 16, 2026).

\bibitem{statcounter_desktop_os}
StatCounter Global Stats, ``Desktop Operating System Market Share Worldwide,'' \url{https://gs.statcounter.com/os-market-share/desktop/worldwide} (accessed June 27, 2026).

\bibitem{atlasv2}
A. Riddle et al., ``ATLASv2: ATLAS Attack Engagements, Version 2,'' arXiv preprint arXiv:2401.01341, 2024.

\bibitem{deeplog}
M. Du et al., ``DeepLog: Anomaly Detection and Diagnosis from System Logs through Deep Learning,'' Proc. ACM CCS, 2017.

\bibitem{nodoze}
W. U. Hassan et al., ``NoDoze: Combatting Threat Alert Fatigue with Automated Provenance Triage,'' Proc. NDSS, 2019.

\bibitem{atlas}
A. Alsaheel et al., ``ATLAS: A Sequence-based Learning Approach for Attack Investigation,'' Proc. USENIX Security, 2021.

\bibitem{airtag}
W. U. Hassan et al., ``AIRTAG: Towards Automated Attack Investigation by Unsupervised Learning with Log Texts,'' Proc. IEEE Symposium on Security and Privacy, 2023.

\bibitem{attack}
MITRE, ``MITRE ATT\&CK,'' \url{https://attack.mitre.org/} (accessed June 16, 2026).

\bibitem{lolbas}
LOLBAS Project, ``Living Off The Land Binaries, Scripts and Libraries,'' \url{https://lolbas-project.github.io/} (accessed June 16, 2026).

\bibitem{react}
S. Yao et al., ``ReAct: Synergizing Reasoning and Acting in Language Models,'' arXiv preprint arXiv:2210.03629, 2022.

\bibitem{clouseau}
Y. H. Lee et al., ``CLOUSEAU: A Hierarchical Multi-Agent Framework for Automated Attack Investigation,'' arXiv preprint arXiv:2505.01799, 2025.
\end{thebibliography}

\end{document}
% END OF MAIN.TEX
